\section{Integral orbits and gradings of a fixed order}\label{sec:rigidity}

An ungraded Morita list does not specify the components of a crossed
order.  This section gives the additional rigidity needed to recover
both their group labels and the identity marking.  There are three
steps: an orbit lemma over $\cO$, a tangent calculation for gradings
of a fixed total order, and a rank argument which passes from Morita
classes to such fixed orders.  The orbit method is related to the
geometry of rigid lattices and Picard groups in
\cite{EiseleGeometry2022}; the grading argument below works for every
finite grading group.

\subsection{Integral orbit finiteness}

\begin{lemma}[Integral orbit finiteness]\label{lem:orbit}
Let $J$ be a smooth affine group scheme of finite type over $\cO$,
acting on an affine scheme $V$ of finite type over $\cO$.  For
$x\in V(\cO)$ let $T_xV$ denote the integral tangent module along
the section $x$.  The points for which the orbit derivative
\[
 \Lie(J)\longrightarrow T_xV
\]
is surjective belong to only finitely many $J(\cO)$-orbits.
\end{lemma}

\begin{proof}
We bound the special-fibre orbits and then cut each one by a fixed
integral slice.  The latter has only finitely many possible generic
points arising from the sections under consideration.

\emph{The integral tangent and horizontal components.}
Embed $V$ as a closed subscheme of $\mathbb A^r_{\cO}$, with
finitely many defining equations.  The module $T_xV$ is the kernel
of their Jacobian map on $\cO^r$.  It is therefore saturated in
$\cO^r$.  Write $u=\rank_{\cO}T_xV$.  Surjectivity of the orbit
derivative and saturation imply that its coordinate matrix has a
unit minor of size $u$.  This follows, for example, from Smith
normal form over the discrete valuation ring $\cO$.

After extension to $K_0$, this derivative has rank $u$, which is
also $\dim_{K_0}T_{x_{K_0}}V_{K_0}$.  Every irreducible component
of $V_{K_0}$ through $x_{K_0}$ consequently has dimension at most
$u$.  Let $V_u$ be the reduced closure in $V$ of the union of the
generic-fibre irreducible components of dimension at most $u$.
There are finitely many such components.  Each of their horizontal
closures has special fibre of dimension at most its generic-fibre
dimension.  This is the dimension theorem for an irreducible
finite-type scheme over a valuation ring
\cite[Lemma~33.19.2, Tag~0B2J]{StacksProject}.  Thus
\begin{equation}\label{eq:fibre-dimension}
 \dim (V_u)_k\leq u.
\end{equation}
The section $x$ factors through $V_u$, because its generic point
does and $\cO\hookrightarrow K_0$ is injective.

\emph{The special-fibre orbits.}
Let $x_0$ be the reduction of $x$.  The unit minor remains nonzero
modulo $p$, so the special-fibre orbit of $x_0$ has dimension at
least $u$.  This orbit lies in $(V_u)_k$.  One way to verify that
last assertion is to lift every element of $J(k)$ to $J(\cO)$,
using smoothness and completeness, and then apply it to the section
$x$.  The transformed generic point still lies on components of
dimension at most $u$.  Equation~\eqref{eq:fibre-dimension} now
shows that the orbit has dimension exactly $u$.

A locally closed orbit of dimension $u$ in a scheme of dimension
at most $u$ contains an open subset of an irreducible component of
dimension $u$.  Two distinct orbits cannot both contain dense open
subsets of the same component.  There are therefore finitely many
possible special-fibre orbits for this $u$.

\emph{A fixed Hensel slice.}
Fix one such orbit and a representative $x_0$.  By lifting elements
of $J(k)$, move all integral points under consideration so that
their reduction equals $x_0$.  Choose $u$ coordinate functions
whose orbit derivative has a nonzero $u$-minor at $x_0$, and fix
integral lifts $a_1,\ldots,a_u$ of their values there.  For each
such integral point $x$, the map
\[
 J\longrightarrow\mathbb A^u_{\cO},\qquad
 g\longmapsto\text{the selected coordinates of }g x
\]
is smooth near the identity: $J$ is smooth and its relative
differential there is surjective.  Hensel lifting supplies
$g\equiv1\pmod p$ for which the selected coordinates of $gx$
are exactly the $a_i$.  Thus every integral orbit has a
representative in the fixed affine slice
\[
 W=V\cap\{\text{selected coordinates }=a_1,\ldots,a_u\}.
\]

For any resulting point $y$, the selected-coordinate differential
is an isomorphism on $T_{y_{K_0}}V_{K_0}$: that space has dimension
$u$, and the same minor is a unit since $y$ reduces to $x_0$.
It follows that $T_{y_{K_0}}W_{K_0}=0$.  Such a point is isolated
in the finite-type $K_0$-scheme $W_{K_0}$.  A noetherian scheme
has only finitely many isolated points.  Each generic point gives
at most one integral section, again by the injection
$\cO\hookrightarrow K_0$.  Hence there are finitely many
representatives in this slice.  Taking the finite union over the
special-fibre orbits and the possible ranks $0\leq u\leq r$
proves the lemma.
\end{proof}

The integral tangent module in this proof need not reduce to the
whole tangent space at $x_0$.  Saturation and the unit minor are
what the proof uses.  In particular, no smoothness assumption on
$V$, nor on an automorphism scheme occurring as $V$, is needed.
The rank-zero case is included by using no slice coordinates.  The
argument uses smoothness of $J$ for lifting and for its differential;
it does not require $J$ to be connected.

\subsection{Block rigidity and labelled gradings}

\begin{lemma}[Hochschild rigidity and outer automorphisms]
\label{lem:HH}
Let $A$ be an $\cO$-order Morita equivalent to a finite direct sum
of finite-group block orders.  Then $\HH^1_{\cO}(A)=0$, and
$\Out_{\cO}(A)$ is finite.
\end{lemma}

\begin{proof}
For a finite group $G$, the conjugation permutation lattice and
Shapiro's lemma give
\[
 \HH^1_{\cO}(\cO G)
 \cong H^1(G,\cO G_{\mathrm{conj}})
 \cong\bigoplus_{[g]}H^1(C_G(g),\cO)=0.
\]
The action on the last coefficient module is trivial, and
$\Hom(C_G(g),\cO_{\mathrm{add}})=0$ because $C_G(g)$ is finite
and $\cO$ is torsion-free.  Hochschild cohomology splits over
finite direct products and is Morita invariant, proving the
first assertion.  Thus every $\cO$-linear derivation of $A$ is
inner.

The unit group scheme $A^\times$ is an open subscheme of the
affine space underlying $A$, defined by invertibility of the
regular-representation determinant.  It is smooth.  The
automorphisms of $A$ form an affine scheme of finite type: impose
the multiplicative and unital equations on an invertible linear
map.  Let $A^\times$ act by postcomposition with inner
automorphisms.  At an automorphism, compose with its inverse to
identify the tangent module with the derivations of $A$.
The orbit derivative then consists of the inner derivations,
so it is surjective.  Lemma~\ref{lem:orbit} says exactly that
there are finitely many cosets modulo inner automorphisms.
\end{proof}

\begin{theorem}[Gradings of a fixed integral order]\label{thm:gradings}
Let $A$ be an $\cO$-order with $\HH^1_{\cO}(A)=0$, and let $H$
be any fixed finite group.  There are only finitely many
$H$-labelled gradings of $A$ up to $\cO$-algebra automorphism.
In fact there are only finitely many orbits under inner
automorphisms.
\end{theorem}

\begin{proof}
We apply Lemma~\ref{lem:orbit} to the space of decompositions
compatible with multiplication.  The key point is that the tangent
calculation uses cancellation in $H$, without averaging over $H$.

A grading $A=\bigoplus_{h\in H}A_h$ is specified by its
orthogonal projections $e_h\in\End_{\cO}(A)$.  These satisfy
\begin{align*}
 e_he_j&=\delta_{h,j}e_h,& \sum_{h\in H}e_h&=1,\\
 e_h(1_A)&=\delta_{h,1}1_A,&
 e_l(e_g(x)e_h(y))&=0\quad(l\ne gh).
\end{align*}
The last equations need only be imposed on a fixed finite basis
of $A$.  They define an affine finite-type grading scheme, on
which $A^\times$ acts by conjugation.

We compute its integral tangent module at a grading.  A
first-order deformation of the direct-sum decomposition is
uniquely represented by an off-diagonal $\cO$-linear map
$\delta:A\to A$, where
\[
 e_h\delta e_h=0\quad(h\in H).
\]
The deformed summands are $(1+\varepsilon\delta)A_h$ over
$\cO[\varepsilon]/(\varepsilon^2)$.  This description follows
by differentiating the equations for orthogonal idempotent
projections, or by writing the summands as graphs over the
original summands.  It uses no division by $|H|$.

For $x\in A_g$ and $y\in A_h$, the linearized multiplicative
equations say that
\begin{equation}\label{eq:grading-derivative}
 \delta(xy)-\delta(x)y-x\delta(y)
\end{equation}
has zero component in every degree other than $gh$.  It also has
zero $gh$-component.  The first term is off degree $gh$ by
definition.  A summand of $\delta(x)y$ can have degree $gh$ only
if its first factor has degree $g$, by cancellation in $H$; that
component of $\delta(x)$ is zero.  The same reasoning applies
to $x\delta(y)$.  Hence Equation~\eqref{eq:grading-derivative}
vanishes, and $\delta$ is a derivation.

The hypothesis gives $\delta=[a,-]$ for some $a\in A$.
Conjugation by $1+\varepsilon a$ produces the specified tangent
to the grading.  Thus the orbit derivative from the smooth
group $A^\times$ is surjective onto every integral tangent
module.  Lemma~\ref{lem:orbit} proves the assertion.
\end{proof}

\begin{remark}\label{rem:fixed-total}
The total order is fixed in Theorem~\ref{thm:gradings}.  Finiteness
of all crossed orders on a fixed identity order is a different
assertion.  For example, for odd $p$ the based $C_p$-crossed orders
\[
 \cO[t]/(t^p-a),\qquad a\in\cO^\times,
\]
have classes parametrized by $\cO^\times/(\cO^\times)^p$.
The principal-unit quotient is infinite: the $p$-adic logarithm
identifies $(1+p\cO)/(1+p\cO)^p$ with
$p\cO/p^2\cO\cong k$.  Their total orders vary.
The crossed-product finiteness theorem in the published
\cite[Corollary~4.9]{EiseleReduction2021} assumes a $p'$-grading
group.  Theorem~\ref{thm:gradings} uses a different hypothesis and
has just been proved also for groups divisible by $p$.
\end{remark}

\subsection{From Morita classes to based graded types}

For a crossed order on a fixed identity order, fixing the grading
group also fixes the total rank.  This converts a finite Morita list
into a finite isomorphism list, to which the grading theorem applies.

\begin{proposition}[Retaining the identity marking]\label{prop:marked}
Fix a basic block order $R$ and a finite group $H$.  Let
$\mathcal A$ be a collection of $H$-crossed orders on $R$, each
Morita equivalent to a finite-group block.  If the total orders
in $\mathcal A$ have finitely many $\cO$-Morita classes, then
$\mathcal A$ has finitely many based graded isomorphism classes.
\end{proposition}

\begin{proof}
\emph{The total order.}
Every order in the collection has rank $|H|\rank_{\cO}R$.
Choose a basic representative in each of its finitely many
Morita classes.  Any order in that class is an endomorphism
order of a projective generator
$\bigoplus_i P_i^{m_i}$, with positive integer multiplicities
$m_i$ and the $P_i$ ranging over the finitely many indecomposable
projectives.  Its endomorphism order contains
$M_{m_i}(\End(P_i))$ as a direct $\cO$-module summand, so its
rank bounds $m_i^2$.  There are consequently only finitely many
ungraded isomorphism types in the collection.

\emph{The grading and the marking.}
Lemma~\ref{lem:HH} and Theorem~\ref{thm:gradings} give finitely
many labelled $H$-gradings on every such total order.  Fix one
graded representative with identity component isomorphic to
$R$.  Two identifications of that component with $R$ differ by
an element of $\Aut_{\cO}(R)$.  Inner changes extend to graded
automorphisms of the total order by conjugation with a unit in
degree one.  Since $\Out_{\cO}(R)$ is finite by
Lemma~\ref{lem:HH}, there are finitely many remaining markings.
These are precisely the based graded types in the assertion.
\end{proof}
